\documentclass[10pt,a4paper]{amsart}
\usepackage[margin=19mm]{geometry}
\usepackage{amsmath,amssymb,mathtools}
\usepackage[scr=rsfs,cal=euler]{mathalfa}
\usepackage{xfrac}
\usepackage[unicode,hidelinks,bookmarks=false]{hyperref}
\newcommand{\ie}{i.e.\ }
\newcommand{\cf}{cf.\ }
\newcommand{\resp}{respectively\ }
\newcommand{\Subsection}{\subsection}
\providecommand{\nobreakdash}{\nobreak\hbox{-}}
\setlength{\emergencystretch}{2em}
\multlinegap0pt
\usepackage{bm}
\usepackage{bbm}
\usepackage{dsfont}
\usepackage{xspace}
\usepackage{cancel}
\usepackage{mathtools}
\usepackage{amsthm}
\makeatletter
\@ifundefined{theorem}{\newtheorem{theorem}{Theorem}}{}
\@ifundefined{lemma}{\newtheorem{lemma}[theorem]{Lemma}}{}
\makeatother
\DeclareMathOperator{\Pol}{Pol} % Polarization
\newcommand{\KL}{\mathcal{L}}
\newcommand{\eps}{\varepsilon} % Epsilon variant #listed
\newcommand{\qtext}[1]{\quad\text{#1}\quad} % Quad spaced text #listed
\newcommand{\ra}{\rightarrow} % Right arrow #listed
\newcommand{\tensor}{\otimes} % Tensor product #listed
\newcommand{\tightlist}{\setlength{\itemsep}{0pt}\setlength{\parskip}{0pt}}
\title[Fa\`a di Bruno is Taylor Composition]{Fa\`a di Bruno is Taylor Composition}
\author{Heinrich Hartmann}
\date{Source: arXiv:2606.26133v2 (CC BY 4.0)}
\begin{document}
\maketitle

\noindent\emph{Source and licence.} Fa\`a di Bruno is Taylor Composition. Version arXiv:2606.26133v2 (CC BY 4.0), distributed under the Creative Commons Attribution 4.0 International licence, as recorded in the arXiv licence metadata for this exact version and shown on the abstract page https://arxiv.org/abs/2606.26133v2. Full rights notice: licenses/arxiv-2606.26133-CC-BY-4.0.txt.\par\smallskip

\noindent\textit{Editorial scope.} Counted unit: the Theorem whose source label is \texttt{thm:fdb-comp} in the article (source0.tex lines 750--762, source label \texttt{thm:fdb-comp}), delivered together with the article's own complete original proof of it (source0.tex lines 764--823, one \texttt{proof} environment of 60 lines). No mathematics is written, repaired or re-derived: statement and proof are the article's own text line for line. Required context retained uncounted: lem:poly-props (lines 574--596); lem:lipschitz (lines 723--730). Every definition, display and label of those lines is retained unchanged. Omitted from this package and not redistributed: 1--573, 597--722, 731--749, 763--763, 824--1240. Nothing inside a retained excerpt is omitted or altered: every retained body is delivered exactly as the article's lines are written, so no omission receipt of any kind accompanies this package. Citations: the retained excerpts carry no citation command at all, so no bibliography entry of the article is transcribed and no reference is redistributed. Reference adaptation: none was needed and none was made; every reference inside the delivered unit resolves inside it, so no unresolved reference or citation is produced. Printed slips kept literal: no printed word, symbol, hypothesis or number of the counted statement or of its proof was altered. Layout and class: the article's own class is \texttt{article} with packages including inputenc, fontenc, lmodern, graphicx, geometry, setspace, enumitem, tocloft, amsmath, amssymb, mathtools, amsthm, longtable, natbib; the wrapper is the lane's pinned \texttt{amsart} editorial wrapper at 10pt on A4, which restates the theorem-like environments the retained text needs and adds the article's own macro definitions that the retained text needs (\texttt{\textbackslash KL}, \texttt{\textbackslash Pol}, \texttt{\textbackslash eps}, \texttt{\textbackslash qtext}, \texttt{\textbackslash ra}, \texttt{\textbackslash tensor}, \texttt{\textbackslash tightlist}), with extra wrapper packages (bm, bbm, dsfont, xspace, cancel, mathtools). The delivered numbering is the unit's own. Assets: no retained span contains a figure environment, a graphics-inclusion command, a PDF-page inclusion, a table or a TikZ picture; no image file is copied into this package, the package declares no asset dependency and no image file is redistributed with it. Licence: version arXiv:2606.26133v2 of this article is distributed under the Creative Commons Attribution 4.0 International licence, as recorded in the arXiv licence metadata for this exact version and shown on the abstract page https://arxiv.org/abs/2606.26133v2. A full rights notice is delivered at licenses/arxiv-2606.26133-CC-BY-4.0.txt. Characters: every retained excerpt is pure ASCII, so the delivered engine, XeLaTeX with its default Latin Modern text font, reports no missing character.
\begin{lemma}[Polynomial Properties]

\label{lem:poly-props} Let \(p : E \ra F\), \(q: F \ra G\) be polynomial
maps.

\begin{enumerate}
\def\labelenumi{\arabic{enumi}.}
\tightlist
\item
  (Composition) The composition \(q \circ p: E \ra G\) is a polynomial
  with \(\deg(q \circ p) \leq \deg(q) \cdot \deg(p)\).
\item
  (Lipschitz) There exists \(\eps > 0, C>0\) such that
  \(\|\Delta(p; 0; x)\| \le C \|x\|\) for all \(\|x\| < \eps\).
\item
  (Upper Vanishing) We have \(\| \pi_{> k}(p)(x)\| / \|x\|^{k} \ra 0\)
  for \(x \ra 0\).
\item
  (Lower Vanishing) If \(\|p(x)\| / \|x\|^{k} \ra 0\) for \(x \ra 0\),
  and \(k \geq \deg(p)\) then \(p = 0\).
\end{enumerate}

\end{lemma}

\begin{lemma}[Lipschitz Continuity]

\label{lem:lipschitz} If \(\phi \in C^k(U_E, F)\) with \(k \geq 1\),
then \(\phi\) is locally Lipschitz continuous: for every \(x \in U_E\)
there exist \(C > 0, \eps > 0\) such that
\(\| \Delta(\phi; x; v) \| \leq C \|v\|\) for all \(\|v\| < \eps\).

\end{lemma}

\begin{theorem}[Taylor Composition]

\label{thm:fdb-comp} Let \(E,F,G\) be Banach spaces. Let
\(U_E \subset E\) be open, \(x \in U_E\), and \(\phi \in C^k(U_E, F)\).
Let \(U_F \subset F\) be open with \(y = \phi(x) \in U_F\) and
\(\psi \in C^k(U_F, G)\).

Then \(\psi \circ \phi \in C^k(U_E, G)\) and the reduced Taylor
polynomials compose: \[
    T_{\ast}^k(\psi \circ \phi;\, x) = \pi_{\leq k}\bigl( T_{\ast}^k(\psi;\, y) \circ T_{\ast}^k(\phi;\, x) \bigr).
    \]

\end{theorem}

\begin{proof}

By translation, we may assume \(x=0\) and \(\phi(x) = y = 0\) as well as
\(\psi(y) = 0\) throughout the proof. In particular
\(\Delta(\phi; 0) = \phi\), \(\Delta(\psi; 0) = \psi\).

Choose Taylor decompositions \(\phi = P + R\) and \(\psi = Q + S\), with
polynomials \(P,Q\) of degree \(\le k\) and remainders \(R,S\) vanishing
to order \(k\) at \(0\). Now write
\(\psi\circ\phi = Q(\phi) + S(\phi)=\pi_{\le k}(Q\circ P) + E\), with:
\[
  E := Q(\phi) + S(\phi) - Q(P) + \pi_{> k}(Q\circ P).
\] We have to show that \(\|E(x)\| / \|x\|^k \to 0\) for \(x \to 0\).

Clearly \(\| \pi_{> k}(Q\circ P) \| / \|x\|^k \to 0\) for \(x \to 0\) by
Lemma \ref{lem:poly-props}.

To see that \(\|S(\phi(x))\|/\|x\|^k \to 0\) we argue as follows: By the
Lipschitz property of \(\phi\) (Lemma \ref{lem:lipschitz}), we find
\(\delta > 0, C > 0\) so that \(\|\phi(x)\| < C \|x\|\) for all
\(\|x\| < \delta\). Since \(S\) is a Taylor residual we have
\(\|S(y)\| < \eta(y) \|y\|^k\) with \(\eta(y) \to 0\) as \(y \to 0\).
Hence \(\|S(\phi(x))\|/\|x\|^k \leq \eta(\phi(x)) C^k \to 0\) as
\(x \to 0\).

For the term \(Q(P+R) - Q(P)\), we decompose
\(Q(y) = \sum_{\ell=0}^k Q_\ell(y) = \sum_\ell \frac{1}{\ell!} B_\ell[y^{\tensor \ell}]\)
where \(Q_\ell = \pi_\ell Q\) and \(B_\ell = \Pol_\ell Q\). Then \[
  Q_\ell(P+R) - Q_\ell(P)
  = \sum_{\substack{i+j=\ell\\ j\ge1}} \frac{1}{i!\,j!}\,B_\ell\big(P^{\tensor i}, R^{\tensor j}\big).
\]

It suffices to show that each summand
\(B_\ell\big(P^{\tensor i}, R^{\tensor j})\) is in \(o(\|x\|^k)\). To
this end we note that:

\begin{itemize}
\tightlist
\item
  \(\|B_\ell(P(x)^{\tensor i}, R(x)^{\tensor j})\| \leq \|B_\ell\| \cdot \|P(x)\|^i \|R(x)\|^j\)
  as \(B_\ell \in \KL^s(^\ell F,G)\) bounded.
\item
  \(\|P(x)\| \le C\|x\|\) for \(\|x\| < \delta\) as \(P\) is Lipschitz
  (Lemma \ref{lem:poly-props}) and \(P(0) = 0\).
\item
  \(\|R(x)\| < \eta(x) \|x\|^k\) with \(\eta(x) \ra 0\) for \(x \to 0\)
  as \(R\) is a Taylor remainder.
\end{itemize}

So that: \[
  \|B_\ell(P^{\tensor i}(x),R^{\tensor j}(x))\| / \|x\|^k
  \leq
  \|B_\ell\| \; C^i \; \eta(x)^j \; \|x\|^{i + (j - 1) k}
  \to 0 \qtext{as} x \to 0
\] For the last step, note that \(j \geq 1\) so the factor
\(\eta(x)^j \to 0\) as \(x \to 0\), and \(i + (j - 1) k \geq 0\) hence
\(\|x\|^{i + (j - 1) k}\) stays bounded (and even goes to zero for
\(j \geq 2\)).

\end{proof}

\end{document}
